% Part: incompleteness
% Chapter: theories-computability
% Section: first-incompleteness

\documentclass[../../../include/open-logic-section]{subfiles}

\begin{document}

\olfileid{inc}{tcp}{inc}

\olsection{$\Th{Q}$-এর কোনো সম্পূর্ণ, সঙ্গতিপূর্ণ,
  !!^{axiomatizable} প্রসারণ নেই}

\begin{thm}
\ollabel{thm:first-incompleteness}
$\Th{Q}$-এর কোনো সম্পূর্ণ, সঙ্গতিপূর্ণ, !!{axiomatizable} প্রসারণ নেই।
\end{thm}

\begin{proof}
আমরা ইতিমধ্যে জানি যে $\Th{Q}$-এর কোনো সঙ্গতিপূর্ণ, নির্ণেয় প্রসারণ
নেই। কিন্তু $\Th{T}$ সম্পূর্ণ এবং !!{axiomatizable} হলে সেটি নির্ণেয়।
% BN-SRC-240 TARGET-CORRECTION-BEGIN
\par\smallskip\noindent\textnormal{\small\textbf{উৎস-সংশোধন (BN-SRC-240):}
হিমায়িত প্রমাণের শেষ বাক্যে লেমাটির অনুমান \texttt{!!\{axiomatizable\}}
থেকে বদলে \texttt{!!\{axiomatized\}} লেখা হয়েছে। উপপাদ্যের ঘোষিত
অনুমান এবং ঠিক আগের লেমার সঙ্গে মিলিয়ে বাংলা প্রমাণে
\texttt{!!\{axiomatizable\}} পুনরুদ্ধার করা হয়েছে।}
% BN-SRC-240 TARGET-CORRECTION-END
\end{proof}

\begin{explain}
এই উপপাদ্যটি গ্যোডেলের ১৯৩১ সালের প্রথম অসম্পূর্ণতা উপপাদ্যের মূল
সূত্রায়ন থেকে খুব দূরে নয়। অপেক্ষাকৃত আধুনিক পরিভাষা বাদ দিলে মূল
পার্থক্যগুলি হলো: গ্যোডেল ``সঙ্গতিপূর্ণ''-এর জায়গায়
``$\omega$-সঙ্গতিপূর্ণ'' লিখেছিলেন; আর গণনসাধ্যতার বিধিবদ্ধ ধারণাটি তখনও
প্রতিষ্ঠিত হয়নি বলে তিনি পূর্ণ সাধারণতায় ``!!{axiomatizable}'' বলতে
পারেননি। (পরের দশকে গণনসাধ্যতার বিধিবদ্ধ মডেলগুলি বিকশিত হয়; গ্যোডেলও
এই কাজে অংশ নেন। যেসব তত্ত্বে গ্যোডেলীয় ঘটনাটি ঘটে, সেগুলির ধরন
চিহ্নিত করাই এই বিকাশের অন্যতম প্রধান উদ্দেশ্য ছিল।)

উপপাদ্যটি বলে, সম্পূর্ণতা, সঙ্গতি ও !!{axiomatizability}---এই তিনটিই এক
সঙ্গে পাওয়া যায় না। তবে যেকোনো একটি ছেড়ে দিলে বাকি দুটি পাওয়া যায়:
$\Th{Q}$ সঙ্গতিপূর্ণ ও গণনসাধ্যভাবে স্বতঃসিদ্ধায়িত, কিন্তু সম্পূর্ণ নয়;
অসঙ্গতিপূর্ণ তত্ত্বটি সম্পূর্ণ এবং গণনসাধ্যভাবে স্বতঃসিদ্ধায়িত (যেমন,
$\{ \eq/[0][0] \}$ দ্বারা), কিন্তু সঙ্গতিপূর্ণ নয়; আর পাটিগণিতের সত্য
বাক্যগুলির সেট সম্পূর্ণ ও সঙ্গতিপূর্ণ, কিন্তু গণনসাধ্যভাবে স্বতঃসিদ্ধায়িত
নয়।
\end{explain}
\end{document}
